\documentclass{article}

\usepackage{microtype}
\usepackage{graphicx}
\usepackage{subcaption}
\usepackage{booktabs}
\usepackage{hyperref}

\newcommand{\theHalgorithm}{\arabic{algorithm}}

\usepackage[accepted]{icml2026}

\usepackage{amsmath}
\usepackage{amssymb}
\usepackage{mathtools}
\usepackage{amsthm}

\usepackage[capitalize,noabbrev]{cleveref}

\usepackage{algorithm}
\usepackage{algorithmic}
\usepackage{listings}
\usepackage{xcolor}

\lstset{
  basicstyle=\ttfamily\scriptsize,
  breaklines=true,
  breakatwhitespace=false,
  frame=single,
  rulecolor=\color{gray!50},
  backgroundcolor=\color{gray!5},
  keywordstyle=\color{blue!70}\bfseries,
  commentstyle=\color{green!50!black},
  stringstyle=\color{orange!80!black},
  numbers=left,
  numberstyle=\tiny\color{gray},
  numbersep=4pt,
  xleftmargin=12pt,
  language=Python,
  tabsize=4,
  showstringspaces=false,
}

\icmltitlerunning{\phantom{x}}

\begin{document}

\twocolumn[
  \icmltitle{SalmoTrace-BERT: Atribusi Sumber Genomik \textit{Salmonella enterica}
    Menggunakan Fitur SNP, Frekuensi K-mer,\\
    dan Embedding DNABERT-2}

  \icmlsetsymbol{equal}{*}

  \begin{icmlauthorlist}
    \icmlauthor{Muhammad Raihaan Perdana}{}
    \icmlauthor{Danendra Shafi Athallah}{}
    \icmlauthor{M. Abizzar Gamadrian}{}
    \icmlauthor{Kenzie Raffa Ardhana}{}
  \end{icmlauthorlist}

  \icmlkeywords{epidemiologi genomik, atribusi sumber, Salmonella,
    DNABERT-2, SNP, k-mer, machine learning, keamanan pangan}

  \vskip 0.3in
]

\printAffiliationsAndNotice{}

\begin{abstract}
Mengidentifikasi asal isolat \textit{Salmonella enterica}, apakah berasal
dari rantai pangan atau bukan, merupakan tantangan utama dalam surveilans
keamanan pangan berbasis genomik. Kami menyajikan SalmoTrace-BERT, sebuah
pipeline Python yang melakukan atribusi sumber biner pada 168 rakit genom
publik dari NCBI Pathogen Detection, menggunakan tiga representasi fitur
genomik komplementer: (i) matriks SNP berbasis referensi (8.157 posisi),
(ii) frekuensi pentanukleotida k-mer bebas-alignment (1.024 fitur), dan
(iii) embedding kontekstual DNABERT-2 (768 dimensi). Klasifikasi dievaluasi
dengan validasi silang \textsc{StratifiedGroupKFold} yang dikelompokkan
berdasarkan kluster SNP NCBI untuk mencegah kebocoran data dari isolat
berkerabat. Studi ablasi menunjukkan bahwa fitur k-mer mencapai macro F1
terbaik 0,655, sementara DNABERT-2 tanpa fine-tuning hanya menghasilkan
F1\,=\,0,515. Perlu dicatat bahwa performa klasifikasi secara keseluruhan
masih terbatas, hal ini disebabkan ukuran dataset yang relatif kecil
(168 isolat) dan homogenitas genomik inherent isolat \textit{Salmonella}
publik yang membatasi daya diskriminasi sinyal genomik.
\end{abstract}

% =========================================================================
\section{Pendahuluan}
% =========================================================================

\textit{Salmonella enterica} adalah penyebab utama penyakit bawaan makanan
di seluruh dunia, bertanggung jawab atas sekitar 93 juta kasus dan 155.000
kematian setiap tahun \cite{who2023salmonella}. Dalam konteks keamanan pangan,
kemampuan melacak asal isolat, apakah dari rantai produksi pangan seperti
ternak dan unggas, atau dari lingkungan non-rantai pangan seperti sampel
klinis, merupakan fondasi investigasi wabah berbasis genomik.

\textit{Whole-genome sequencing} (WGS) telah mentransformasi epidemiologi
dengan menyediakan matriks \textit{single-nucleotide polymorphism} (SNP) yang
menangkap variasi genetik pada resolusi tinggi \cite{loman2015rapid}. Alat
pemanggil SNP berbasis referensi seperti Snippy \cite{illingworth2021snippy}
dan BLAST+ \cite{camacho2009blast} telah menjadi standar investigasi wabah,
di mana jarak Hamming berpasangan dari matriks core-SNP digunakan untuk
pengelompokan hierarkis dan atribusi sumber \cite{achtman2012mlst,nguyen2020source}.
Namun, pipeline alignment penuh bersifat mahal secara komputasi, dan daya
diskriminasinya bergantung pada keragaman genomik kumpulan isolat.

Sebagai pelengkap, profil frekuensi k-mer bebas-alignment, seperti yang
digunakan pada sketsa MinHash \cite{ondov2016mash}, dapat menangkap sinyal
komposisi nukleotida yang relevan tanpa alignment eksplisit. Di sisi lain,
model bahasa DNA pra-latih DNABERT \cite{ji2021dnabert} dan DNABERT-2
\cite{zhou2023dnabert2} menghasilkan embedding kontekstual 768 dimensi yang
secara teoritis mengkodekan motif sekuens relevan-sumber; namun keduanya
memerlukan fine-tuning pada tugas genomik hilir untuk mencapai performa
kompetitif.

Perbandingan sistematis ketiga jenis representasi ini, SNP, k-mer, dan
embedding, untuk atribusi sumber \textit{Salmonella} dengan evaluasi yang
ketat terhadap kebocoran data dari isolat berkerabat belum pernah dilaporkan.
Penelitian ini mengisi kesenjangan tersebut. Perlu ditekankan bahwa proyek
ini merupakan \textbf{simulasi forensik genomik} menggunakan 168 isolat data
publik NCBI; ukuran dataset yang terbatas dan homogenitas genomik inherent
isolat publik menyebabkan performa klasifikasi belum optimal, hal ini menjadi
batasan eksplisit yang dibahas pada Bagian~\ref{sec:hasil}.

\textbf{Rumusan masalah:}
\begin{enumerate}
  \item Apakah frekuensi k-mer bebas-alignment dapat menyamai atau melampaui
        fitur SNP untuk atribusi sumber?
  \item Apakah DNABERT-2 tanpa fine-tuning memberikan informasi komplementer
        terhadap fitur SNP?
  \item Seberapa diskriminatif pengelompokan hierarkis berbasis SNP untuk
        memisahkan isolat rantai pangan dari non-rantai pangan?
\end{enumerate}

\textbf{Kontribusi:}
\begin{itemize}
  \item Pipeline Python end-to-end (\textit{SalmoTrace-BERT}) mencakup QC
        genom, ekstraksi SNP, komputasi k-mer, embedding DNABERT-2, dan
        klasifikasi multi-fitur dengan metadata AMR.
  \item Protokol evaluasi group-aware (StratifiedGroupKFold per kluster SNP
        NCBI) yang memitigasi kebocoran dari isolat berkerabat.
  \item Studi ablasi komparatif pada 168 genom \textit{S.~enterica} dengan
        interpretasi biologis dan analisis keterbatasan eksplisit.
\end{itemize}

% =========================================================================
\section{Metode}
% =========================================================================

\subsection{Dataset dan Kontrol Kualitas}

Kami mengunduh 168 rakit genom \textit{S.~enterica} (format FASTA) beserta
metadata dari NCBI Pathogen Detection \cite{ncbi_pathogen}. Genom referensi
adalah \textit{S.~enterica} serovar Typhimurium LT2 (aksesi GCF\_000006945.2).
Setiap rakit dikenai tiga filter kualitas: ukuran genom 4--6\,Mb, fraksi-N
$\leq$5\%, dan kontig $\leq$500. Seluruh 168 isolat lulus QC.

Label sumber isolasi dinormalisasi dari 23 kategori mentah menjadi target
biner seimbang: \texttt{food\_chain\_associated} (FCA) dan
\texttt{non\_food\_chain\_associated} (NFCA), masing-masing 84 isolat.
Cakupan geografis mencakup 13 negara; 48\% isolat berasal dari Amerika Serikat.

\subsection{Ekstraksi Fitur}

\textbf{Matriks SNP.}
Pemanggilan SNP berbasis referensi dilakukan dengan pencocokan k-mer (k\,=\,21),
menghasilkan 8.157 posisi variabel. Nukleotida dikodekan sebagai bilangan bulat
(A$\to$0, C$\to$1, G$\to$2, T$\to$3) membentuk matriks $168\times8.157$.

\textbf{Fitur k-mer.}
Frekuensi pentanukleotida (5-mer) dihitung genome-wide, menghasilkan matriks
$168\times1.024$ yang ternormalisasi dalam [0,\,1].

\textbf{Embedding DNABERT-2.}
Sekuens dipotong menjadi jendela 512\,bp (maks.\ 500 per isolat, stride 256\,bp),
ditokenisasi, lalu dilewatkan melalui DNABERT-2-117M \cite{zhou2023dnabert2}.
Embedding per-isolat diperoleh via \textit{masked mean pooling} atas seluruh
jendela. Model digunakan sebagai ekstraktor beku tanpa fine-tuning.

\textbf{Fitur AMR.}
Vektor biner kehadiran gen resistansi antibiotik dari metadata NCBI disertakan
pada kondisi ablasi tertentu.

\textbf{Attention MIL.}
Model \textit{Multiple Instance Learning} berbasis atensi dilatih atas embedding
per-jendela DNABERT-2, memungkinkan pembobotan jendela secara diskriminatif.

\subsection{Klasifikasi dan Evaluasi}

Semua klasifikasi menggunakan Random Forest (200 estimator,
\texttt{random\_state}\,=\,42) kecuali dinyatakan lain. Protokol utama adalah
\textsc{StratifiedGroupKFold} 5-lipat dengan \texttt{snp\_cluster} NCBI sebagai
grup, mencegah isolat berkerabat muncul di lipatan pelatihan dan pengujian.
Metrik: F1 makro, akurasi seimbang, F1 berbobot, dan silhouette score.

\subsection{Gambaran Umum Pipeline}
\label{sec:pipeline}

Gambar~\ref{fig:flowchart} mengilustrasikan alur pipeline secara keseluruhan,
dari input data hingga laporan forensik.

\begin{figure*}[ht]
  \vskip 0.1in
  \raggedright
  \includegraphics[height=0.60\textheight, keepaspectratio]{pipeline_flowchart}
  \caption{Alur pipeline SalmoTrace-BERT: dari input FASTA dan metadata hingga
    atribusi sumber forensik melalui tiga jalur ekstraksi fitur (SNP, k-mer,
    dan DNABERT-2).}
  \label{fig:flowchart}
  \vskip -0.1in
\end{figure*}

\subsection{Cuplikan Kode Utama}

Berikut tiga fungsi yang menjadi inti pipeline. Kode lengkap tersedia di
repositori GitHub proyek.

\textbf{(1) Ekstraksi Matriks SNP} (\texttt{src/snp/extractor.py}).
Fungsi \texttt{build\_core\_snp\_matrix()} membangun matriks SNP berbasis
referensi menggunakan indeks k-mer vectorised dan diagonal chaining per kontig:

\begin{lstlisting}[caption={Ekstraksi matriks core-SNP berbasis referensi.}]
def build_core_snp_matrix(
    genomes, reference_seq, k=21,
    min_core_fraction=0.95,
    min_isolate_callability=0.10,
):
    ref_seq_up = reference_seq.upper()
    ref_len    = len(ref_seq_up)
    # Bangun indeks k-mer unik dari referensi (vectorised NumPy)
    ref_hashes, ref_pos = _build_ref_kmer_index(ref_seq_up, k)
    # Alokasi matriks uint8: n_isolat x ref_len (A/T/G/C/N)
    matrix = np.zeros((len(genomes), ref_len), dtype=np.uint8)
    for i, (acc, contigs) in enumerate(genomes.items()):
        _fill_isolate_row(contigs, ref_hashes,
                          ref_pos, matrix[i], k)
    # Filter: core positions + remove invariant sites
    callable_frac = (matrix > 0).mean(axis=0)
    core_mask = callable_frac >= min_core_fraction
    snp_matrix = matrix[:, core_mask]
    variant_mask = snp_matrix.max(0) != snp_matrix.min(0)
    return pd.DataFrame(snp_matrix[:, variant_mask],
                        index=list(genomes.keys()))
\end{lstlisting}

\textbf{(2) Komputasi K-mer} (\texttt{src/data/kmer.py}).
Fungsi \texttt{extract\_kmer\_features()} menghitung frekuensi 5-mer yang
ternormalisasi secara genome-wide tanpa alignment:

\begin{lstlisting}[caption={Komputasi frekuensi pentanukleotida (5-mer).}]
def extract_kmer_features(genomes, k=5):
    kmers = _all_kmers(k)   # 4^k = 1024 kombinasi
    rows  = {}
    for acc, seq in genomes.items():
        seq_up = seq.upper()
        counts = Counter()
        for i in range(len(seq_up) - k + 1):
            window = seq_up[i: i + k]
            if all(c in "ATGC" for c in window):
                counts[window] += 1
        total = sum(counts.values()) or 1
        # Normalisasi: frekuensi dalam [0, 1]
        rows[acc] = {km: counts.get(km, 0) / total
                     for km in kmers}
    return pd.DataFrame(rows).T
\end{lstlisting}

\textbf{(3) Embedding DNABERT-2} (\texttt{src/embedding/encoder.py}).
Fungsi \texttt{embed\_sequence()} menghasilkan vektor 768 dimensi per jendela
dengan \textit{masked mean pooling} atas \textit{hidden state} terakhir:

\begin{lstlisting}[caption={Encoding satu jendela DNA dengan DNABERT-2.}]
def embed_sequence(seq, tokenizer, model, device="cpu"):
    inputs = tokenizer(seq, return_tensors="pt",
                       truncation=True, max_length=512)
    input_ids     = inputs["input_ids"].to(device)
    attention_mask = inputs["attention_mask"].to(device)
    with torch.no_grad():
        outputs = model(input_ids)
    hidden = outputs[0]  # (1, seq_len, 768)
    # Masked mean pooling: abaikan token padding
    mask = attention_mask.unsqueeze(-1).float()
    emb  = (hidden * mask).sum(dim=1) \
           / mask.sum(dim=1).clamp(min=1e-9)
    return emb.squeeze(0).cpu().numpy()  # (768,)
\end{lstlisting}

% =========================================================================
\section{Hasil dan Diskusi}
\label{sec:hasil}
% =========================================================================

\subsection{Analisis Data Eksploratif (EDA)}

Sebelum klasifikasi, kami melakukan analisis eksploratif untuk memahami
karakteristik dataset. Tabel~\ref{tab:eda} merangkum statistik utama.

\begin{table}[ht]
  \caption{Statistik deskriptif dataset SalmoTrace-BERT.}
  \label{tab:eda}
  \begin{center}
    \begin{small}
      \begin{tabular}{lr}
        \toprule
        Karakteristik                        & Nilai         \\
        \midrule
        Jumlah isolat total                  & 168           \\
        Isolat FCA / NFCA                    & 84 / 84       \\
        Jumlah negara asal                   & 13            \\
        Proporsi isolat dari USA             & 48\%          \\
        Posisi SNP teridentifikasi           & 8.157         \\
        Jarak SNP minimum antar-isolat       & $<$10         \\
        Jarak SNP maksimum antar-isolat      & 993           \\
        Rata-rata jarak SNP                  & 505,85        \\
        Jumlah kluster hierarkis (k=5)       & 5             \\
        Kluster SNP NCBI unik                & 52            \\
        Dimensi embedding DNABERT-2          & 768           \\
        Fitur k-mer (5-mer)                  & 1.024         \\
        \bottomrule
      \end{tabular}
    \end{small}
  \end{center}
  \vskip -0.1in
\end{table}

Distribusi kelas biner seimbang sempurna (84:84), menghilangkan bias
mayoritas-kelas. Namun, distribusi geografis sangat tidak merata (48\% dari
USA), sehingga generalisasi ke populasi global terbatas. Dari 23 kategori
sumber isolasi mentah, 14 kategori memiliki kurang dari 3 sampel, indikasi
ketidakseimbangan kelas halus yang dapat mendestabilisasi estimasi F1
per-kelas. Rentang jarak SNP yang luas (min\,$<$10 hingga maks\,993)
mengindikasikan isolat dengan kekerabatan genetik sangat dekat di satu sisi
dan isolat yang jauh berkerabat di sisi lain; hal ini juga menjelaskan
mengapa rasio separasi kluster ($1{,}6\times$) tergolong rendah.

\subsection{Pengelompokan Hierarkis Berbasis SNP}

Pengelompokan hierarkis (average linkage) menghasilkan lima kluster
(silhouette 0,284, dipilih berdasarkan inspeksi dendrogram). Rata-rata jarak
SNP intra-kluster 443,1 vs.\ 711,0 antar-kluster menghasilkan rasio separasi
$1{,}6\times$ yang rendah, mencerminkan rentang jarak yang sangat lebar
(min\,$<$10, maks\,993): pasangan klonal dan isolat jauh-berkerabat hadir
bersamaan. Kluster pertama didominasi NFCA, sementara kluster tiga dan empat
mencakup campuran FCA/NFCA hampir setara.

Tabel~\ref{tab:clustering} menunjukkan ARI dan NMI terhadap empat kolom
metadata. ARI\,=\,0,001 terhadap \texttt{snp\_cluster} NCBI (52 kluster,
resolusi $\approx$50 SNP) mencerminkan perbedaan granularitas: lima kluster
kami merepresentasikan filogeni tingkat atas, bukan pengelompokan sumber
spesifik. NMI\,=\,0,209 untuk \texttt{isolation\_source} mengonfirmasi sinyal
bersama yang lemah.

\begin{table}[t]
  \caption{Validasi kluster terhadap metadata (ARI / NMI). Pengelompokan
    dengan hierarchical clustering average-linkage, $k\!=\!5$.}
  \label{tab:clustering}
  \begin{center}
    \begin{small}
      \begin{tabular}{lrrr}
        \toprule
        Kolom metadata       & ARI      & NMI    & \#Kat \\
        \midrule
        isolation\_source    & 0,025    & 0,209  & 22    \\
        serovar              & $-$0,032 & 0,074  & 3     \\
        snp\_cluster (NCBI)  & 0,001    & 0,191  & 52    \\
        geo\_loc\_name       & $-$0,020 & 0,161  & 38    \\
        \bottomrule
      \end{tabular}
    \end{small}
  \end{center}
  \vskip -0.1in
\end{table}

Meski tidak memulihkan label sumber, lima kluster ini berperan kritis sebagai
\textit{stratifikasi grup} pada StratifiedGroupKFold: isolat klonal dalam
kluster yang sama dipisahkan ke lipatan berbeda, mencegah kebocoran data
(lihat Bagian~\ref{sec:hasil}).

\subsection{Studi Ablasi Atribusi Sumber}

Tabel~\ref{tab:ablation} merangkum seluruh 14 konfigurasi fitur-model yang
dievaluasi. Semua konfigurasi melampaui baseline dummy (F1\,=\,0,282) secara
substansial, namun performa tertinggi (F1\,=\,0,655) masih terbilang moderat,
konsisten dengan keterbatasan dataset yang kecil dan homogenitas genomik
dataset yang dibahas sebelumnya. Berikut analisis per kelompok eksperimen.

\begin{table}[t]
  \caption{Perbandingan lengkap seluruh konfigurasi fitur dan model,
    StratifiedGroupKFold 5-lipat. \textbf{Cetak tebal} = terbaik.
    RF=RandomForest, LR=Logistic Regression, SVC=Support Vector,
    MIL=Attention-MIL.}
  \label{tab:ablation}
  \begin{center}
    \begin{scriptsize}
      \setlength{\tabcolsep}{4pt}
      \begin{tabular}{llccc}
        \toprule
        Fitur            & Model & F1$_\text{mac}$ & Bal.Acc & F1$_w$ \\
        \midrule
        \multicolumn{5}{l}{\textit{Berbasis SNP}} \\
        SNP              & RF    & 0,634$\pm$0,10  & 0,666   & 0,657  \\
        SNP              & LR    & 0,578$\pm$0,13  & 0,645   & 0,591  \\
        SNP              & SVC   & 0,589$\pm$0,11  & 0,618   & 0,621  \\
        \midrule
        \multicolumn{5}{l}{\textit{Berbasis DNABERT-2}} \\
        DNABERT-2        & RF    & 0,515$\pm$0,09  & 0,563   & 0,523  \\
        DNABERT-2        & LR    & 0,485$\pm$0,04  & 0,518   & 0,502  \\
        DNABERT-2        & SVC   & 0,472$\pm$0,03  & 0,499   & 0,493  \\
        DNABERT-2        & MIL   & 0,622           & 0,637   & 0,642  \\
        \midrule
        \multicolumn{5}{l}{\textit{Berbasis K-mer}} \\
        K-mer            & RF    & \textbf{0,655}$\pm$0,09 & \textbf{0,676} & \textbf{0,674} \\
        \midrule
        \multicolumn{5}{l}{\textit{Gabungan (Hibrida)}} \\
        SNP + DNABERT    & RF    & 0,635$\pm$0,10  & 0,665   & 0,660  \\
        K-mer + AMR      & LR    & 0,589$\pm$0,06  & 0,604   & 0,616  \\
        SNP + K-mer + AMR& RF    & 0,586$\pm$0,08  & 0,639   & 0,604  \\
        \midrule
        \multicolumn{5}{l}{\textit{Baseline}} \\
        AMR only         & LR    & 0,413$\pm$0,06  & 0,504   & 0,456  \\
        Dummy            & ---   & 0,282$\pm$0,08  & 0,500   & 0,252  \\
        \bottomrule
      \end{tabular}
    \end{scriptsize}
  \end{center}
  \vskip -0.1in
\end{table}

\textbf{K-mer unggul.}
Profil frekuensi pentanukleotida mencapai macro F1 terbaik 0,655 (Bal.Acc
0,676), melampaui SNP meski tanpa alignment. Komposisi nukleotida global
menangkap bias penggunaan kodon dan preferensi dinukleotida skala-genom yang
berkaitan dengan adaptasi FCA secara holistik, sekaligus menjadi pilihan
paling ringan secara komputasi. Standar deviasi $\pm$0,09 konsisten di semua
konfigurasi RF, mencerminkan variansi dari ukuran lipatan ($\approx$33 isolat)
bukan instabilitas model.

\textbf{SNP kompetitif.}
SNP-RF mencapai F1\,=\,0,634 (hanya 0,021 di bawah k-mer). Skor silhouette
SNP (0,284) jauh lebih tinggi dari k-mer (0,103), namun keunggulan geometris
ini tidak berkonversi menjadi klasifikasi lebih baik di bawah evaluasi
group-aware karena isolat klonal dipisahkan oleh stratifikasi grup. RF
mengungguli SVC (0,589) dan LR (0,578), mengonfirmasi struktur non-linear di
ruang 8.157 dimensi.

\textbf{DNABERT-2 sebagai representasi perantara.}
DNABERT-2 berfungsi sebagai lapisan pemrosesan representasi, bukan fitur
klasifikasi langsung. Embedding beku 768 dimensi belum disesuaikan untuk
sinyal sumber \textit{Salmonella} sehingga RF langsung hanya menghasilkan
F1\,=\,0,515; model linear lebih buruk lagi (LR 0,485; SVC 0,472),
mengindikasikan struktur yang tidak terorganisir linear. Attention MIL yang
mengagreagsi embedding per-jendela (512\,bp) secara selektif meningkatkan
F1 ke 0,622 (+0,107), menunjukkan bahwa informasi sumber terkonsentrasi di
jendela genomik tertentu yang dapat diidentifikasi mekanisme atensi.

\textbf{Hibrida dan AMR.}
Penggabungan fitur tidak memberikan peningkatan: SNP+DNABERT (0,635) hampir
identik SNP saja; K-mer+AMR (0,589) dan tiga fitur (0,586) justru lebih
rendah dari K-mer tunggal karena fitur AMR memperkenalkan noise. AMR saja
(F1\,=\,0,413) mengonfirmasi gen resistansi antibiotik tidak membawa sinyal
sumber yang memadai. RF secara konsisten mengungguli LR dan SVC di semua
kelompok fitur.

\subsection{Analisis Kebocoran Data}

Untuk mengkuantifikasi dampak pemisahan group-aware, kami membandingkan
StratifiedGroupKFold group-aware dengan pemisahan terStratifikasi naif.
Peningkatan F1 maksimum di bawah pemisahan naif adalah
$\Delta\text{F1}_{\max} = +0{,}088$, mengonfirmasi inflasi moderat ketika
isolat berkerabat bocor antar-lipatan. Semua hasil yang dilaporkan menggunakan
protokol group-aware.

\subsection{Interpretasi Biologis}

Keunggulan fitur k-mer atas alignment berbasis SNP memiliki penjelasan biologis
yang masuk akal: isolat yang berasosiasi dengan rantai pangan mungkin berbagi
bias penggunaan kodon, pola komposisi GC, atau preferensi dinukleotida skala
genom yang berkaitan dengan adaptasi terhadap matriks pangan kaya nutrien dan
lingkungan saluran cerna ternak. Sinyal komposisi ini ditangkap secara holistik
oleh profil 5-mer, tetapi terdilusi dalam matriks SNP posisional yang jarang
di mana alel spesifik pada masing-masing dari 8.157 posisi harus secara
individual memprediksi sumber.

Nilai ARI yang sangat rendah (0,001) antara lima kluster hierarkis kami dan
52 kluster SNP yang ditetapkan NCBI menunjukkan bahwa pengelompokan kami
menangkap tingkat granularitas genomik yang berbeda. Kluster SNP NCBI
didefinisikan pada resolusi lebih halus; menggabungkannya menjadi lima grup
memulihkan sinyal geografis atau filogeni luas, bukan sinyal spesifik-sumber.

% =========================================================================
\section{Kesimpulan}
% =========================================================================

Kami telah menyajikan SalmoTrace-BERT, sebuah pipeline multi-fitur untuk
atribusi sumber biner \textit{Salmonella enterica} dari data WGS publik.
Temuan utama kami adalah:

\begin{itemize}
  \item Profil k-mer pentanukleotida bebas-alignment (F1\,=\,0,655) sedikit
        mengungguli fitur SNP berbasis alignment (F1\,=\,0,634) untuk
        klasifikasi rantai pangan vs.\ non-rantai pangan, menunjukkan bahwa
        komposisi nukleotida global mengkodekan informasi relevan-sumber.
  \item DNABERT-2 berperan sebagai lapisan pemrosesan representasi perantara,
        bukan fitur klasifikasi primer: embedding beku 768 dimensi dari model
        pra-latih mengkodekan informasi kontekstual multi-spesies yang belum
        disesuaikan untuk sinyal sumber \textit{Salmonella}, sehingga skor RF
        langsung hanya mencapai F1\,=\,0,515. Mekanisme Attention MIL yang
        mengagreagsi embedding per-jendela secara selektif lebih efektif
        memanfaatkan representasi tersebut (F1\,=\,0,622), dan fine-tuning
        tugas-spesifik diharapkan meningkatkan kualitas representasi lebih lanjut.
  \item Evaluasi group-aware sangat penting: pemisahan naif menginflaasi F1
        hingga 0,088 akibat kebocoran keterkaitan.
  \item Rasio separasi genomik yang rendah ($1{,}6\times$) dan ARI mendekati
        nol terhadap kluster NCBI mencerminkan homogenitas genomik rakit
        \textit{Salmonella} yang tersedia secara publik, yang membatasi
        performa klasifikasi absolut.
\end{itemize}

\textbf{Arah penelitian selanjutnya.}
Fine-tuning DNABERT-2 pada data atribusi sumber \textit{Salmonella} berlabel
diharapkan dapat meningkatkan kualitas embedding secara substansial. Perluasan
dataset melampaui 168 isolat, penggabungan metadata wabah epidemiologi nyata,
dan validasi pada koleksi yang beragam secara geografis juga merupakan langkah
penting ke depan. Integrasi dengan alignment core-SNP berbasis Snippy dapat
memberikan posisi SNP yang lebih presisi dibandingkan pendekatan pencocokan
k-mer yang kami gunakan.

% =========================================================================
\section*{Pernyataan Dampak}
% =========================================================================

Penelitian ini bertujuan untuk memajukan metode komputasi dalam surveilans
keamanan pangan berbasis genomik. Dengan menunjukkan bahwa fitur k-mer
bebas-alignment dapat menyaingi metode berbasis SNP untuk atribusi sumber,
pipeline kami menurunkan hambatan komputasi untuk epidemiologi genomik di
lingkungan dengan sumber daya terbatas. Semua data yang digunakan tersedia
secara publik melalui NCBI; tidak ada informasi pasien pribadi yang diakses.
Pipeline ini tidak dimaksudkan untuk digunakan dalam investigasi wabah resmi
tanpa validasi epidemiologi independen, sebagaimana dicatat secara eksplisit
dalam dokumentasi keterbatasan kami.

% =========================================================================
\section*{Perangkat Lunak dan Data}
% =========================================================================

Kode dan dokumentasi pipeline tersedia di GitHub. Dataset diperoleh dari
NCBI Pathogen Detection (\url{https://www.ncbi.nlm.nih.gov/pathogens/}).
Checkpoint model DNABERT-2 tersedia secara publik di Hugging Face
(\texttt{zhihan1996/DNABERT-2-117M}).

% =========================================================================

\bibliography{salmotrace_paper}
\bibliographystyle{icml2026}

\end{document}
